%% 13장 컴퓨터 구조
\chapter{컴퓨터 구조}
\label{ch:architecture}

\abstract{같은 수의 연산을 하는 두 프로그램의 속도가 열 배 넘게 차이 나는 일은 드물지 않다. 연산의 수만 세는 알고리즘 분석으로는 이 차이를 설명할 수 없고, 연산이 실제로 실행되는 기계의 구조를 알아야 한다. 이 장에서는 현대 프로세서의 성능을 좌우하는 네 가지를 차례로 본다. 데이터를 가져오는 비용이 연산의 비용보다 훨씬 큰 기억 계층과 캐시, 연산 하나의 지연 시간과 단위 시간당 처리량을 구별하게 하는 파이프라인, 조건 분기의 결과를 미리 짐작하는 분기 예측, 한 명령으로 여러 값을 처리하는 SIMD이다. 마지막으로 프로그램이 보는 주소와 실제 메모리를 잇는 가상 메모리를 간단히 다룬다. 이 장의 수치는 모두 Apple M4 Pro에서 직접 작성한 작은 측정 프로그램으로 잰 것이고, 기계마다 다르지만 그 비율과 원리는 거의 모든 현대 프로세서에서 같다.}

\section{기억 계층}
\label{sec:memory-hierarchy}

\subsection{지연 시간의 사다리}

프로세서는 덧셈 하나를 $1$나노초보다 훨씬 짧은 시간에 끝내지만, 주기억장치(DRAM)에서 값 하나를 가져오는 데는 그 백 배가 넘는 시간이 걸린다. 이 격차는 수십 년 동안 벌어지기만 했는데\cite{wulf1995}, 프로세서는 이를 메우려고 작고 빠른 기억장치를 여러 층으로 둔다. 레지스터 바로 다음에 \term{L1 캐시}{L1 cache}가 있고, 그다음에 더 크고 느린 L2 캐시가 있으며, 그 뒤에 DRAM이 있다. 최근에 쓴 데이터와 그 근처의 데이터를 가까운 층에 복사해 두면 대부분의 접근을 빠른 층에서 처리할 수 있다는 것이 캐시의 원리이다.

그 사다리의 높이를 재어 보자. 무작위 순서로 연결된 목록을 따라가면서 각 원소에서 다음 원소의 위치를 읽게 하면, 다음 주소를 알아야 다음 읽기를 시작할 수 있으므로 읽기 하나하나의 지연 시간이 그대로 드러난다. 이 \term{포인터 추적}{pointer chasing}을 목록 전체의 크기를 바꿔 가며 M4 Pro의 성능 코어에서 잰 결과가 \cref{tab:latency}이다. 목록이 $96$KiB 이하이면 읽기 하나에 약 $0.9$ns이고, $512$KiB에서 $4$MiB 사이는 $6$--$7$ns, $64$MiB 이상은 약 $115$ns이다. 이 기계의 성능 코어 하나에는 $128$KiB의 L1 데이터 캐시가 있고 성능 코어들이 $16$MiB의 L2 캐시를 나누어 쓰므로, 세 계단은 각각 L1, L2, DRAM에 해당한다. L1과 DRAM 사이의 차이는 $100$배를 넘는다.

\begin{table}[t]
\begin{sdtable}{colspec={lccccc}, row{1}={font=\small\bfseries}}
\toprule
목록의 크기 & $96$KiB & $512$KiB & $4$MiB & $12$MiB & $256$MiB \\
\midrule
읽기 하나의 지연 & $0.89$ns & $5.93$ns & $7.21$ns & $12.1$ns & $122$ns \\
\bottomrule
\end{sdtable}
\caption{포인터 추적으로 잰 Apple M4 Pro 성능 코어의 읽기 지연. 각 원소는 캐시 줄 하나를 차지하고 무작위 순서로 이어진다. $12$MiB는 L2 안이지만 주소 변환의 비용(\cref{sec:virtual-memory})이 더해진다.}
\label{tab:latency}
\end{table}

\subsection{캐시 줄과 지역성}

캐시는 바이트 하나씩이 아니라 \term{캐시 줄}{cache line}이라는 고정 크기의 덩어리로 데이터를 옮긴다. M4 Pro의 캐시 줄은 $128$바이트이므로, binary64 값 하나를 읽으면 그 이웃 $15$개가 함께 올라온다. 그래서 메모리를 차례로 훑는 프로그램은 캐시 줄 하나를 가져오는 비용을 $16$개의 값이 나누어 낸다. 이처럼 최근에 쓴 주소 근처를 곧 다시 쓰는 성질을 \term{공간 지역성}{spatial locality}, 같은 주소를 곧 다시 쓰는 성질을 \term{시간 지역성}{temporal locality}이라 하며, 캐시는 두 지역성이 있는 프로그램에서만 효과를 낸다. 게다가 프로세서는 차례로 진행되는 접근 패턴을 알아채고 다음 캐시 줄을 미리 가져오는 \term{선인출}{prefetching}까지 하므로, 순차 접근은 지연 시간을 거의 감춘다.

$4096\times4096$ 행렬의 원소를 행 순서로 더하면 원소 하나에 $0.48$ns가 걸리지만, 열 순서로 더하면 $2.7$ns가 걸린다. 행렬은 행이 연속해서 저장되므로(17장), 열 순서로 훑으면 원소마다 $32$KiB씩 건너뛰어 매번 다른 캐시 줄을 읽게 되기 때문이다. 연산의 수는 똑같은데 속도는 약 $5.7$배 차이가 난다. 행렬 곱셈처럼 같은 데이터를 여러 번 쓰는 계산에서는 행렬을 캐시에 들어가는 작은 블록으로 나누어 한 블록을 다 쓰고 나서 다음 블록으로 가는 \term{블록화}{blocking, tiling}가 핵심 기법이 되며, 고성능 BLAS의 행렬 곱은 L1, L2, 주기억장치의 층마다 블록의 크기를 맞추어 데이터가 각 층에 머무는 동안 최대한 많이 쓴다\cite{lam1991blocked,goto2008}.

\begin{instardust}[\code{docs/design/array.md} §3.3]
배열의 원소별 연산은 입력과 출력의 메모리 배치를 보고 반복의 순서를 정한다. 모든 피연산자가 연속이면 하나의 평평한 반복으로 처리하고, 그렇지 않으면 축들을 보폭의 크기에 따라 다시 배열해 가장 안쪽 반복이 가장 촘촘한 축을 돌게 한다. 그래서 행 순서 배열과 열 순서 배열을 섞어 쓰면 어느 쪽 순서로도 접근이 연속이 되지 않아 블록 단위로 옮기는 느린 길로 떨어지는데, 설계 문서는 이 경우를 따로 짚어 둔다. 배열의 메모리 배치와 보폭은 17장에서 자세히 다룬다.
\end{instardust}

\section{파이프라인}
\label{sec:pipeline}

\subsection{지연 시간과 처리량}

프로세서는 명령 하나를 여러 단계로 나누어 처리하고, 각 단계를 맡은 회로가 동시에 서로 다른 명령의 서로 다른 단계를 처리한다. 공장의 조립 라인과 같은 이 구조를 \term{파이프라인}{pipeline}이라 한다. 그 결과 연산 하나에는 두 가지 시간이 생긴다. 하나는 입력이 들어가서 결과가 나오기까지의 \term{지연 시간}{latency}이고, 다른 하나는 새 연산을 얼마나 자주 시작할 수 있는가, 곧 \term{처리량}{throughput}이다. 현대 프로세서에는 같은 종류의 연산 장치가 여러 개 있어서 처리량은 지연 시간보다 훨씬 높다.

차이를 재어 보자. FMA $x\gets\operatorname{fma}(x,a,b)$를 같은 변수에 대해 되풀이하면 각 연산이 앞 연산의 결과를 기다려야 하므로 지연 시간이 드러나고, M4 Pro에서 연산 하나에 약 $0.78$ns가 걸린다. 서로 독립인 $16$개의 변수에 대해 번갈아 같은 연산을 하면 연산 하나에 약 $0.059$ns, 곧 $13$배 넘게 빠르다. 연산 장치가 쉬지 않으려면 지연 시간 동안 다른 독립 연산이 그 자리를 채워야 하는데, 이 관계는 대기열 이론의 Little 법칙\cite{little1961}과 같은 꼴이다.
\[
\text{동시에 진행 중이어야 할 연산의 수}=\text{지연 시간}\times\text{처리량}.
\]
지연 시간이 $4$사이클이고 매 사이클 $4$개의 연산을 시작할 수 있다면 서로 독립인 연산이 $16$개쯤 진행 중이어야 장치가 놀지 않는다.

\subsection{의존 사슬을 끊기}

이 원리를 합에 적용해 보자. 9장의 재귀 합 $s\gets s+x_i$는 덧셈마다 앞 덧셈의 결과를 기다리는 하나의 긴 의존 사슬이다. $2^{20}$개의 binary64 값을 더할 때 누산기 하나로는 원소 하나에 $0.48$ns가 걸리지만, 누산기 $16$개를 두고 원소를 번갈아 나누어 넣은 뒤 마지막에 모으면 $0.064$ns로 $7$배 넘게 빨라진다. 덧셈의 수는 같고 의존 관계만 바뀌었을 뿐이다. 9장에서 보았듯이 이 재배치는 오차의 상한도 줄인다. 합의 순서를 나무 모양으로 바꾸는 것이 정확도와 속도 모두에 이롭다.

\begin{instardust}[\code{csrc/base/pairwise.hpp}]
stardust의 쌍별 합에서 잎의 누산기를 $16$개로 둔 이유가 바로 이것이다. 주석에 따르면 $16$개의 독립된 사슬을 컴파일러가 NEON 레지스터 여덟 개(레지스터 하나에 binary64 둘)로 채운다. numpy의 같은 알고리즘은 누산기 $8$개와 길이 $128$의 블록을 쓰는데, 잎을 넓힌 stardust의 판은 한 스레드에서 긴 합이 $1.2$--$1.6$배 빨랐다고 설계 문서에 적혀 있다. 대신 합하는 순서가 달라지므로 마지막 비트가 numpy와 다를 수 있는데, 이것은 8장에서 본 결정성의 약속, 곧 같은 기계에서는 언제나 같은 비트를 낸다는 약속과는 충돌하지 않는다.
\end{instardust}

\subsection{순서를 바꾸어 실행하기}

현대 프로세서는 명령을 프로그램에 적힌 순서대로가 아니라 입력이 준비된 순서대로 실행하고, 결과만 원래 순서대로 확정한다. 이 \term{비순차 실행}{out-of-order execution}은 Tomasulo가 1967년에 IBM System/360 Model 91을 위해 고안한 방법에서 시작되었다\cite{tomasulo1967}. 비순차 실행 덕분에 프로그래머가 의존 사슬을 명시적으로 끊지 않아도 서로 독립인 연산들이 겹쳐 실행되지만, 프로세서가 내다볼 수 있는 명령의 범위는 수백 개로 한정되어 있다. 하나의 긴 의존 사슬, 예컨대 수백만 번 이어지는 재귀 합은 프로세서가 아무리 멀리 내다보아도 겹쳐 실행할 수 있는 것이 없다. 그래서 위의 $16$개 누산기처럼 알고리즘 수준에서 독립성을 만들어 주어야 한다.

\section{분기 예측}
\label{sec:branch-prediction}

파이프라인에는 조건 분기가 골칫거리이다. 분기의 조건을 계산하기 전에 다음 명령을 파이프라인에 넣어야 하는데, 어느 쪽으로 갈지 모르기 때문이다. 프로세서는 과거의 분기 결과에서 다음 결과를 짐작하는 \term{분기 예측기}{branch predictor}를 두고 그 짐작대로 미리 실행한다\cite{smith1981branch}. 짐작이 맞으면 비용이 거의 없지만, 틀리면 미리 실행한 일을 모두 버리고 다시 시작해야 하므로 십수 사이클을 잃는다.

예측기는 규칙적인 패턴을 잘 맞히고 무작위한 패턴은 맞히지 못한다. $0$부터 $255$까지의 무작위 정수 $2^{22}$개에서 $128$ 이상인 것만 더하는 반복을 분기로 쓰면, 원소 하나에 무작위 순서의 배열에서는 $2.5$ns가 걸리고 같은 배열을 정렬한 뒤에는 $0.22$ns가 걸린다. 정렬된 배열에서는 분기의 결과가 처음 절반은 모두 거짓, 나머지 절반은 모두 참이어서 예측기가 거의 언제나 맞히기 때문이다. 같은 데이터, 같은 연산인데 $11$배 차이가 난다.

분기 예측이 실패할 수 있는 곳에서는 분기를 아예 없애는 것이 좋다. 조건에 따라 두 값 가운데 하나를 고르는 일은 비교의 결과를 마스크로 만들어 산술로 처리할 수 있고(ARM의 \code{CSEL}과 NEON의 \code{BSL} 같은 명령이 이것을 한다), 그러면 데이터에 상관없이 늘 같은 시간이 걸린다. 이분 탐색이 좋은 예이다. 정렬된 배열에서 값의 자리를 찾는 이분 탐색은 비교마다 왼쪽과 오른쪽 가운데 하나로 가는데, 찾는 값이 무작위이면 이 분기는 절반의 확률로 틀린다. 분기 없이 쓰면 매 단계가 ``비교 결과에 따라 시작점에 반 길이를 더하거나 더하지 않는'' 산술이 된다.

\begin{instardust}[\code{csrc/base/interval.hpp}]
보간 영역과 배열의 \code{interp}가 함께 쓰는 구간 탐색 \code{count\_le\_lanes}는 분기 없는 이분 탐색을 $L$개의 점에 대해 나란히 한다. 모든 점이 같은 길이로 반씩 줄여 가므로 반복의 모양이 같고, 각 점의 탐색이 서로 독립이어서 캐시를 놓친 읽기 $L$개가 동시에 진행된다. 앞의 Little 법칙을 메모리에 적용한 것으로, 읽기 하나의 지연 시간은 줄일 수 없지만 여러 읽기를 겹치면 처리량은 늘릴 수 있다. 점들이 정렬되어 있으면 직전 점의 결과에서 출발하는 지수 탐색을 써서 비교의 수 자체를 줄인다.
\end{instardust}

\section{SIMD}
\label{sec:simd}

\subsection{한 명령으로 여러 값을}

같은 연산을 여러 데이터에 적용하는 계산은 \term{SIMD}{single instruction, multiple data} 명령으로 한꺼번에 할 수 있다. ARM의 NEON은 $128$비트 레지스터를 쓰므로 binary64 둘, $32$비트 정수 넷, 바이트 열여섯을 한 명령으로 처리한다. 프로세서에 이런 장치가 여러 개 있으므로 SIMD는 파이프라인의 처리량에 곱해지는 또 하나의 병렬성이다.

컴파일러는 단순한 반복을 스스로 SIMD 명령으로 바꾸는데(\term{자동 벡터화}{auto-vectorization}), 이것이 가능하려면 몇 가지 조건이 필요하다. 반복 사이에 의존 관계가 없어야 하고, 반복 안에 예측하기 어려운 분기가 없어야 하며, 입력과 출력의 메모리가 겹치지 않는다는 것을 컴파일러가 알아야 한다. 부동소수점 합처럼 순서를 바꾸면 결과가 달라지는 리덕션은 컴파일러가 마음대로 여러 갈래로 나눌 수 없으므로(8장), 프로그래머가 갈래를 명시적으로 나누어 적어야 한다. 앞 절의 $16$개 누산기가 그런 예이다.

\subsection{자료의 배치}

SIMD는 데이터가 레지스터에 담기 좋은 모양으로 놓여 있을 때 가장 효과적이다. 복소수 배열이 실수부와 허수부를 번갈아 저장하거나 double-double 배열이 높은 부분과 낮은 부분을 번갈아 저장하는 것처럼 서로 다른 성분이 섞여 있으면, 한 레지스터에 같은 성분끼리 모아야 같은 연산을 적용할 수 있다. NEON에는 번갈아 놓인 데이터를 읽으면서 성분별로 갈라 두 레지스터에 담는 \code{LD2} 같은 명령이 있어서 이 일을 따로 하지 않아도 된다.

\begin{instardust}[\code{csrc/htest/rank\_exact.hpp}]
7장에서 본 동점이 있을 때의 동적 계획은 double-double 배열의 한 행에 다른 행을 더하는 일을 수없이 되풀이한다. \code{dd\_axpy}는 \code{vld2q\_f64}로 double-double 두 개를 읽으면서 높은 부분 둘과 낮은 부분 둘을 각각 한 레지스터에 모으고, 10장의 TwoSum과 정규화를 레지스터 단위로 한 뒤 \code{vst2q\_f64}로 다시 번갈아 쓴다. 칸마다의 연산은 스칼라 판과 똑같은 IEEE 연산을 같은 순서로 하므로, 벡터 길과 스칼라 길이 같은 비트를 낸다.
\end{instardust}

\section{가상 메모리}
\label{sec:virtual-memory}

프로그램이 쓰는 주소는 실제 메모리의 주소가 아니라 운영체제가 프로세스마다 만들어 주는 \term{가상 주소}{virtual address}이다. 가상 주소 공간은 \term{페이지}{page}라는 고정 크기의 덩어리로 나뉘고, 운영체제가 관리하는 페이지 표가 각 가상 페이지를 실제 메모리의 페이지에 대응시킨다. macOS의 Apple Silicon에서 페이지의 크기는 $16$KiB이다. 이 구조 덕분에 프로세스끼리 서로의 메모리를 볼 수 없고, 실제로 쓰는 페이지만 실제 메모리를 차지한다.

대가는 주소 변환이다. 메모리를 읽을 때마다 페이지 표를 찾아보면 너무 느리므로, 프로세서는 최근의 변환 결과를 \term{TLB}{translation lookaside buffer}라는 작은 캐시에 둔다. 작업 집합이 TLB가 담을 수 있는 페이지의 범위를 넘으면 주소 변환 자체가 캐시를 놓치고, 데이터 캐시에 값이 있어도 읽기가 느려진다. \cref{tab:latency}에서 $12$MiB의 목록이 L2 캐시 안에 있는데도 $4$MiB보다 느렸던 것이 이 효과이다. 프로그램이 쓰는 메모리가 시간에 따라 어떻게 바뀌는지를 \term{작업 집합}{working set}의 개념으로 처음 정리한 것은 Denning이다\cite{denning1968}.

새로 할당한 메모리는 처음 쓸 때 실제 페이지가 붙는다. 운영체제가 그 페이지를 $0$으로 채워 프로세스에 넘기는 \term{페이지 부재}{page fault} 처리를 하므로, 큰 배열을 처음 쓰는 반복은 같은 반복을 두 번째로 할 때보다 느리다. 성능을 잴 때 측정을 여러 번 되풀이해서 가장 빠른 값을 보는 이유 가운데 하나가 이것이다.

\begin{instardust}[\code{csrc/htest/moments.hpp}]
가설검정 영역의 기술통계 커널에는 할당의 비용에 대한 주석이 있다. 짧은 줄을 수없이 처리할 때 줄마다 작은 벡터를 할당하고 해제하면 계산보다 할당의 비용이 커지는데, 주석은 macOS의 메모리 해제가 시계를 읽기 때문에 한 쌍에 약 $100$ns가 든다고 적는다. 그래서 작은 버퍼는 스택에 두고, 큰 버퍼는 스레드마다 하나씩 두어 다시 쓰며, $2^{16}$개를 넘게 자란 것만 줄이 끝날 때 돌려준다.
\end{instardust}

\section{맺음}

현대 프로세서에서 성능을 정하는 것은 연산의 수보다 데이터의 이동과 의존 관계이다. 데이터를 가까운 층에 두고(지역성), 서로 독립인 일을 충분히 만들어 주고(의존 사슬 끊기), 예측하기 어려운 분기를 산술로 바꾸고, 같은 연산을 여러 값에 한꺼번에 적용하면, 같은 알고리즘이 몇 배에서 수십 배 빨라진다. 이 장의 측정은 한 기계의 것이지만 그 비율은 놀라울 만큼 보편적이다. 다음 장에서는 프로세서 코어 여러 개를 함께 쓰는 병렬 계산으로 넘어가고, 그때 결과의 비트가 스레드의 수와 실행 순서에 따라 달라지지 않게 하는 방법을 다룬다.

컴퓨터 구조 전반은 Hennessy와 Patterson의 교과서\cite{hennessy2019}가 표준이고, 메모리 계층이 프로그램의 성능에 미치는 영향은 Drepper의 글\cite{drepper2007}이 실용적으로 자세하다.

\chapterreferences
